% Etapa 6B -- memoria breve para integracion
% Datos heredados de 3C y auditoria 6A; no se han repetido simulaciones.
\subsection{Verificacion preliminar de uniones y montajes}
Las reacciones exigidas al eje transversal por el modelo biempotrado de la etapa 3C alcanzan una resultante transversal maxima de $7{,}632\ \mathrm{kN}$ por extremo y un momento de empotramiento maximo de $255{,}04\ \mathrm{N\,m}$. Estos maximos no son necesariamente simultaneos y representan solicitaciones de una restriccion ideal, no capacidades validadas de la fijacion.
El eje incorpora roscas M24$\times$2, con area resistente nominal a traccion de aproximadamente $384{,}4\ \mathrm{mm^2}$. Las reacciones transversales no constituyen automaticamente cargas axiales de tuerca. El calculo de precarga, arrancamiento de filetes y presion bajo arandela requiere un corte del montaje, engrane efectivo, materiales y mecanismo de bloqueo.
En cada rodamiento exterior 6206--2RS, una franja axial de $1{,}6\ \mathrm{mm}$ del anillo interior queda sobre un alivio de diametro $29{,}2\ \mathrm{mm}$ en lugar del asiento nominal de $30\ \mathrm{mm}$. Se recomienda revisar la continuidad del asiento y sus hombros antes de aceptar los valores de vida ya calculados.
Para la matriz se propone chaveta paralela de seccion nominal $14\times9\ \mathrm{mm}$, hombro de apoyo y retencion axial positiva independiente. El torque maximo de $88{,}924\ \mathrm{N\,m}$ sobre diametro de $50\ \mathrm{mm}$ equivale a una fuerza tangencial de $3557\ \mathrm{N}$. Con longitud activa hipotetica de $22\ \mathrm{mm}$, el cortante nominal medio seria $11{,}55\ \mathrm{MPa}$ y la presion nominal sobre el cubo $46{,}19\ \mathrm{MPa}$, considerando altura de flanco de $3{,}5\ \mathrm{mm}$. Estos valores no constituyen una verificacion final de la chaveta ni del chavetero real.
Los apoyos del arbol principal todavia carecen de posiciones y tipo definitivos. Se propone una arquitectura localizadora axial y otro apoyo radial no localizador, condicionada a su diseno. Se conserva el empuje maximo de compactacion de $14{,}071\ \mathrm{kN}$ y la resultante radial de correa de hasta $1{,}170\ \mathrm{kN}$ obtenida bajo las hipotesis de 5F. Quedan pendientes el reparto axial, las reacciones radiales, la vida de rodamientos, los anclajes y la capacidad del bastidor.
\paragraph{Alcance.} Las conclusiones anteriores son de ingenieria basica y no autorizan fabricacion sin confirmar geometria, materiales y contactos.